\emph{Preliminaries.} Write $\tau = \langle\cdot\rangle_{0}$ and $u \cdot u' = \tfrac12(uu' + u'u)$ for vectors. By Proposition \claimref{tower.ladder-pairs-satisfy-car-vectors}, $W$ and $W'$ are isotropic and $2v_{i} \cdot w_{j} = \delta_{ij}$, so the form pairs $W$ with $W'$ perfectly. An orthogonal linear map $M$ of $V$ extends uniquely to an automorphism and (composed with the reverse) to an anti-automorphism of the algebra; the adjoint of a split is the anti-automorphism extending the reflection $M_{s}$ that is $+1$ on $V_{+}$ and $-1$ on $V_{-}$, and for the stored split it is the blade inverse of Theorem \claimref{synthesis.chsh-bound-for-two-subsystems-of-a-clifford-algebra-by-a-rational-certificate} (it fixes $e_{j}$ and negates $f_{j}$). The composition of two adjoints is the automorphism extending the product of the two reflections.

\emph{(i).} Let $s$ be adapted and $M = M_{s}$. For $v, v' \in W$, $g(v,v') = 2v \cdot Mv'$; it is symmetric because $M$ is an orthogonal involution ($v \cdot Mv' = Mv \cdot M^{2}v' = Mv \cdot v'$). For $v \ne 0$ the vector $v + Mv$ is fixed by $M$, so it lies in $V_{+}$, and it is nonzero because $W \cap W' = 0$; hence $0 < Q(v + Mv) = 2v \cdot Mv = g(v,v)$, using $Q(v) = Q(Mv) = 0$. Conversely let $g$ be positive definite and symmetric on $W$, and let $\theta : W \to W'$ be the linear map with $2v \cdot \theta v' = g(v,v')$ (it exists and is injective because the pairing is perfect and $g$ nondegenerate). Put $V_{\pm} = \{v \pm \theta v\}$. Then $Q(v \pm \theta v) = \pm g(v,v)$, $(v + \theta v) \cdot (v' - \theta v') = \tfrac12(-g(v,v') + g(v',v)) = 0$, and $V_{+} \cap V_{-} = 0$, so by dimension $V = V_{+} \oplus V_{-}$ is a split; its reflection sends $v$ to $\theta v$, so it is adapted and its form is $g$. The two constructions are inverse to each other. For $A \in \mathrm{GL}_{\ell}(k)$: $2v'_{i} \cdot w'_{j} = \sum_{a}A_{ai}(A^{-1})_{ja} = \delta_{ij}$, so the $e'_{j}$, $f'_{j}$ are orthonormal with squares $1$, $-1$, the reflection of this frame exchanges $v'_{j}$ and $w'_{j}$, and the split is adapted. For $\ell = 1$ a form on $W = kv_{1}$ is $g(v_{1},v_{1}) = c > 0$, with $\theta v_{1} = cw_{1}$ and $Q(v_{1} \pm cw_{1}) = \pm c$; a vector of square $1$ in $V_{+}$ exists exactly when $c$ is a square in $k$. That the stored pairs fix the stored split is $e_{j} = v_{j} + w_{j}$, $f_{j} = v_{j} - w_{j}$. For $A = (t)$: $e' = tv_{1} + t^{-1}w_{1} = \tfrac12(t + t^{-1})e_{1} + \tfrac12(t - t^{-1})f_{1}$, which at $t = 4$ is $\tfrac{17}{8}e_{1} + \tfrac{15}{8}f_{1}$, and $f'e' = (tv_{1} - t^{-1}w_{1})(tv_{1} + t^{-1}w_{1}) = v_{1}w_{1} - w_{1}v_{1} = f_{1}e_{1}$. The space of symmetric forms on $W$ has dimension $\ell(\ell+1)/2$.

\emph{(ii).} If $v'_{j} = \sum_{i}A_{ij}v_{i}$ is another basis of $W$ and $w'_{j}$ the dual basis of $W'$, then $w'_{j} = \sum_{i}(A^{-1})_{ji}w_{i}$ and $\sum_{j}w'_{j}v'_{j} = \sum_{i,a}\bigl(\sum_{j}(A^{-1})_{ji}A_{aj}\bigr)w_{i}v_{a} = \sum_{i}w_{i}v_{i}$. For an adapted split with form $g$ choose a basis $v'_{j}$ of $W$ and put $w'_{j} = v_{j}'^{\dagger_{s}} \in W'$; if $g$ has an orthonormal basis this gives dual bases and $N^{\dagger_{s}} = \sum_{j}(w'_{j}v'_{j})^{\dagger_{s}} = \sum_{j}w'_{j}v'_{j} = N$. In general let $G_{ij} = g(v'_{i},v'_{j})$; the dual basis of $W'$ is $\sum_{i}(G^{-1})_{ji}w'_{i}$, so $N = \sum_{i,j}(G^{-1})_{ji}w'_{i}v'_{j}$ and $N^{\dagger_{s}} = \sum_{i,j}(G^{-1})_{ji}w'_{j}v'_{i} = N$ because $G$ is symmetric. The elements $n_{j} = w_{j}v_{j}$ are commuting idempotents ($w_{j}v_{j}w_{j}v_{j} = w_{j}(1 - w_{j}v_{j})v_{j} = w_{j}v_{j}$), so $N$ satisfies $\prod_{m=0}^{\ell}(N - m) = 0$, the $Q_{n}$ are the idempotents of its spectral decomposition, polynomials in $N$, hence self-adjoint for every adapted split; $Q_{n}$ is the sum of the $\binom{\ell}{n}$ products $\prod_{j \in I}n_{j}\prod_{j \notin I}(1 - n_{j})$, $|I| = n$, each of scalar part $2^{-\ell}$ since $n_{j} = \tfrac12(1 - p_{j})$; and the total parity is $\prod_{j}(1 - 2n_{j}) = \sum_{n}(-1)^{n}Q_{n}$. $Q_{0} = \prod_{j}(1 - n_{j}) = \prod_{j}\tfrac12(1 + p_{j})$ is the vacuum projector of Theorem \claimref{synthesis.born-weights-under-positivity-a-state-sharp-on-an-atom-is-its-trace-functional}, clause (i), with corner $kQ_{0}$; the corner does not depend on the adjoint, so $Q_{0}$, being self-adjoint for every adapted split, is an atom for each, and clause (iii) of that theorem, whose proof uses only that the adjoint is an anti-automorphism with positive definite $\tau(x^{\dagger_{s}}x)$ preserving $\tau$, gives the uniqueness of the sharp state for each.

\emph{(iii).} The condition is linear in $\omega$. Sufficiency: $h(N)$ is self-adjoint for every adapted split by (ii), and $\tau(x^{\dagger_{s}}) = \tau(x)$ for every split (an anti-automorphism extending an orthogonal map preserves the scalar part), so $\tau(h(N)x^{\dagger_{s}}) = \tau((xh(N))^{\dagger_{s}}) = \tau(h(N)x)$. The functionals $\tau(N^{j}\,\cdot)$, $0 \le j \le \ell$, are independent because $1, N, \dots, N^{\ell}$ are and $\tau(y^{\dagger}y) > 0$ for $y \ne 0$. Necessity. Step 1: let $\omega$ be real for the stored split and for the split with $A = \operatorname{diag}(t_{1}, \dots, t_{\ell})$. Then $\omega$ is invariant under the composition of the two adjoints, the automorphism with $v_{j} \mapsto t_{j}^{2}v_{j}$, $w_{j} \mapsto t_{j}^{-2}w_{j}$. The ordered monomials $w_{I}v_{J}$ ($I$, $J$ subsets of the modes) are a basis, and this automorphism multiplies $w_{I}v_{J}$ by $\prod_{j \in J}t_{j}^{2}\prod_{i \in I}t_{i}^{-2}$; taking $t_{j} = 2$ for one $j$ in the symmetric difference of $I$ and $J$ and $t = 1$ elsewhere gives $\omega(w_{I}v_{J}) = 0$ unless $I = J$. So $\omega$ is determined by the numbers $c_{I} = \omega(\prod_{j \in I}n_{j})$; this already proves the statement about splits adapted to each mode separately (dimension $2^{\ell}$; sufficiency there as above, the $n_{j}$ being self-adjoint for such splits). Step 2, $\ell = 2$: the script computes the space of functionals real for the four splits $A = 1$, $\operatorname{diag}(2,1)$, $\operatorname{diag}(1,3)$ and the shear $v'_{2} = v_{1} + v_{2}$; it has dimension $3$ and equals the span of $\tau(N^{j}\,\cdot)$, $j = 0, 1, 2$. The conditions have rational coefficients, so the dimension is the same over every $k$. In particular $c_{\{1\}} = c_{\{2\}}$. Step 3, general $\ell$: fix two modes $i \ne j$ and a set $I$ of other modes, and put $\omega_{I}(y) = \omega(\prod_{m \in I}n_{m}\,y)$ for $y$ in the subalgebra of the modes $i$, $j$. An adapted split of that two-mode algebra, extended by the stored split on the other modes, is an adapted split of the whole, its adjoint fixes $\prod_{m \in I}n_{m}$, which commutes with $y$; so $\omega_{I}$ is real for every adapted split of the two modes, and by Step 2 $\omega_{I}(n_{i}) = \omega_{I}(n_{j})$, that is $c_{I \cup \{i\}} = c_{I \cup \{j\}}$. Hence $c_{I}$ depends only on $|I|$, the space has dimension at most $\ell + 1$, and with sufficiency it is the span of the $\tau(N^{j}\,\cdot)$; the script confirms dimension $4$ for $\ell = 3$. States, in the sense of Premise \claimref{synthesis.premise-positivity-of-values} read for the split $s$: write $h(N) = \sum_{n}h(n)Q_{n}$. If $\omega$ is a state for one adapted split then $\omega(Q_{n}) = h(n)\tau(Q_{n}) \ge 0$, as $Q_{n} = Q_{n}Q_{n}^{\dagger_{s}}$, and $\sum_{n}h(n)\tau(Q_{n}) = 1$; conversely for every adapted split each $x \mapsto \tau(Q_{n}x)/\tau(Q_{n})$ is a state (the functional of an element $Q_{n}Q_{n}^{\dagger_{s}}$; clause (iii) of the CHSH theorem, whose proof applies to the blade inverse of any orthonormal frame, and for a split without orthonormal frame over $k$ by extending $k$), so the convex combinations are. The vacuum state is $q_{0} = 1$; $\prod_{j}\tfrac12(1 + rp_{j}) = \prod_{j}\bigl(\tfrac{1+r}{2}(1 - n_{j}) + \tfrac{1-r}{2}n_{j}\bigr)$ is a function of the $n_{j}$ symmetric in them with coefficient depending only on the number of occupied modes, hence a function of $N$; with distinct $r_{j}$ it is a function of the $n_{j}$.

\emph{(iv).} A state of the stored split that is a state for every adapted split is real for all of them, hence of the form (iii). For the pair state of Theorem \claimref{synthesis.mode-split-of-the-clifford-algebra-and-its-klein-twist-to-the-tensor-product}, clause (v), $\rho_{12}$ contains $cs(w_{1}w_{2} + v_{2}v_{1})$, so its functional $x \mapsto 4\tau(\rho_{12}x)$ is nonzero on the monomial $w_{1}w_{2}$ paired against $v_{2}v_{1}$, a monomial with $I \ne J$, and Step 1 excludes it; the two values are the script's. For the CHSH element: as in the cited clause (vii), $a_{i}b_{j} = x_{i}x'_{j}\,e_{2}f_{1} + x_{i}z'_{j}\,e_{1}p_{2} + z_{i}x'_{j}\,p_{1}e_{2}p_{1} + z_{i}z'_{j}\,p_{1}p_{2}$, using $e_{1}e_{2}p_{1} = e_{2}f_{1}$. The second and third terms are odd, and $e_{2}f_{1} = (v_{2} + w_{2})(v_{1} - w_{1})$ is a combination of monomials $w_{I}v_{J}$ with $I \ne J$ (after ordering, $v_{2}v_{1}$, $w_{1}v_{2}$, $w_{2}v_{1}$, $w_{2}w_{1}$); a functional of the family (iii) vanishes on all of these by Step 1. So $\omega(C) = Z\omega(p_{1}p_{2})$. The family (iii) and the elements $N$, $Q_{n}$ do not depend on the dual bases, so the argument applies verbatim to the directions built from the frame $e'_{j}$, $f'_{j}$ of any adapted split. That a boost changes the positive structure in the Hermitian setting is clause (iv) of Theorem \claimref{synthesis.positive-structure-is-twistor-form-at-imaginary-time}, cited for comparison only. $|\omega(p_{1}p_{2})| = |q_{0} - q_{1} + q_{2}| \le 1$, and $|Z| \le |z'_{1} + z'_{2}| + |z'_{1} - z'_{2}| \le 2$ since $|z_{i}| \le 1$ and $|z'_{j}| \le 1$.
